%%=============================================================================
%% Samenvatting
%%=============================================================================

% TODO: De "abstract" of samenvatting is een kernachtige (~ 1 blz. voor een
% thesis) synthese van het document.
%
% Een goede abstract biedt een kernachtig antwoord op volgende vragen:
%
% 1. Waarover gaat de bachelorproef?
% 2. Waarom heb je er over geschreven?
% 3. Hoe heb je het onderzoek uitgevoerd?
% 4. Wat waren de resultaten? Wat blijkt uit je onderzoek?
% 5. Wat betekenen je resultaten? Wat is de relevantie voor het werkveld?
%
% Daarom bestaat een abstract uit volgende componenten:
%
% - inleiding + kaderen thema
% - probleemstelling
% - (centrale) onderzoeksvraag
% - onderzoeksdoelstelling
% - methodologie
% - resultaten (beperk tot de belangrijkste, relevant voor de onderzoeksvraag)
% - conclusies, aanbevelingen, beperkingen
%
% LET OP! Een samenvatting is GEEN voorwoord!

%%---------- Nederlandse samenvatting -----------------------------------------
%
% TODO: Als je je bachelorproef in het Engels schrijft, moet je eerst een
% Nederlandse samenvatting invoegen. Haal daarvoor onderstaande code uit
% commentaar.
% Wie zijn bachelorproef in het Nederlands schrijft, kan dit negeren, de inhoud
% wordt niet in het document ingevoegd.

\IfLanguageName{english}{%
\selectlanguage{dutch}
\chapter*{Samenvatting}
\selectlanguage{english}
}{}

%%---------- Samenvatting -----------------------------------------------------
% De samenvatting in de hoofdtaal van het document
\chapter*{\IfLanguageName{dutch}{Samenvatting}{Abstract}}

Studenten met Attention Deficit Hyperactivity Disorder (ADHD) ervaren vaak moeilijkheden met
planning, organisatie en tijdsbeheer. Deze problemen hangen samen met beperkingen in executieve
functies zoals werkgeheugen, inhibitie en tijdsperceptie. In het hoger onderwijs, waar studenten
zelfstandig opdrachten moeten plannen en deadlines moeten opvolgen, kan dit leiden tot verhoogde
stress en uitstelgedrag. Hoewel er veel digitale planningstools bestaan, zijn deze meestal ontworpen
voor een brede doelgroep en houden ze onvoldoende rekening met de specifieke noden van studenten
met ADHD.

Deze bachelorproef onderzoekt hoe een mobiele applicatie op maat van studenten met ADHD kan bijdragen
aan beter taakbeheer en tijdsplanning. De centrale onderzoeksvraag luidt: hoe kan een applicatie op maat
van studenten met ADHD bijdragen aan beter taakbeheer en tijdsplanning, met minder stress en uitstelgedrag?
Om deze vraag te beantwoorden werd een literatuurstudie uitgevoerd naar ADHD, executieve functies,
cognitieve belasting en ontwerpprincipes rond cognitieve toegankelijkheid.

Op basis van deze theoretische inzichten werd een proof-of-concept ontwikkeld als mobielgericht prototype.
De applicatie werd geïmplementeerd met React, TypeScript en Tailwind CSS en bevat functionaliteiten voor
het invoeren van taken, het plannen per dag, en een beloningssysteem met punten en streaks. Daarnaast bevat
de applicatie een profielpagina en instellingenpagina. Het prototype werkt volledig client-side en bevat geen backend.

De resultaten van deze bachelorproef tonen aan dat een prikkelarme interface met beperkte keuze-opties,
duidelijke structuur en directe feedback een veelbelovende aanpak vormt om studenten met ADHD te ondersteunen
in hun dagelijkse planning. Door het ontbreken van formele gebruikerstesten kunnen echter geen definitieve conclusies
worden getrokken over effectiviteit in de praktijk. Toekomstig onderzoek kan zich richten op evaluatie met de doelgroep
en uitbreiding van de applicatie met automatische taakopsplitsing en persistente dataopslag.



